\section{Security Considerations}
\label{sec:security}

\subsection{Proposer-trust surface}

The carried-payload model concentrates the irreversible venue
interaction in the proposer: the proposer performs the import-credit and
the debit-export once at \texttt{BuildBlock}, and validators replay the
\emph{results}. This is the same trust surface the order-submission work
established for fills, and it is bounded the same way.

\paragraph{What the proposer can do.} A proposer can build a block whose
carried custody records reflect venue relays it performed, and then
\emph{fail to win the consensus poll} for that block. If the block does
not gain consensus, its custody transitions are not committed to the
canonical chain, yet the proposer already mutated the venue ledger (and,
for a withdraw, may have produced an export UTXO). This is the same
``built-and-relayed but lost the poll'' surface as for fills.

\paragraph{What the proposer cannot do.} The surface is
\emph{conservation-bounded}. A proposer cannot inflate supply:

\begin{itemize}[leftmargin=2em,nosep]
\item It cannot carry a \texttt{credited} larger than the consumed-UTXO
  value --- the import leg refuses, so no such block is even valid
  (Section~\ref{sec:conservation}, fact~2).
\item It cannot carry a \texttt{realized} larger than the account's
  \texttt{available} at relay, nor export more than \texttt{realized}
  --- the export leg refuses.
\item It cannot double-credit a deposit or double-export a withdraw ---
  the consumed-UTXO set and \texttt{seen:<txID>} forbid it
  (Section~\ref{sec:replay}), and both are state the root commits to, so
  a divergent replay outcome is a divergent root.
\end{itemize}

Therefore the worst a malicious or faulty proposer achieves is a
\emph{stranded} custody transition (Section~\ref{sec:strand}), never
created or destroyed value. The blast radius of a single bad block is a
single deposit or a single withdraw escrow, bounded by that
transaction's own value. This matches the accepted single-operator-venue
posture today; the trustless path of Section~\ref{sec:reserved} removes
even the stranding risk by requiring the venue's signature over the
carried results.

\subsection{Reorg}

A reorg that drops an accepted block drops its carried custody records
with it. Two cases:

\begin{itemize}[leftmargin=2em,nosep]
\item \textbf{Deposit in a reorged-out block.} The import consumed a
  source UTXO. If the block is reorged out, the consumed-UTXO set entry
  is rolled back along with the rest of the block's state, so the UTXO
  is spendable again and the deposit can be re-proposed in the new
  canonical chain. No double-credit can occur because the credit only
  exists where the consuming block exists.
\item \textbf{Withdraw in a reorged-out block.} The debit of
  \texttt{available} is rolled back with the block. The export UTXO,
  however, is produced into cross-chain shared memory at
  \texttt{Accept}; a reorg that removes the producing block must also
  remove its shared-memory effect, or the system would have exported
  value with no backing debit. The withdraw path therefore MUST tie the
  export's shared-memory production to the same atomic commit point as
  the block's acceptance, so that reorg removes both or neither.
\end{itemize}

\paragraph{Requirement.} The custody records and their cross-chain
shared-memory effects MUST be committed at a single atomic point with
the block's acceptance --- the same all-or-nothing discipline the proxy
already uses for imports/exports. The carried record is the
in-consensus half; the shared-memory apply is the cross-chain half; they
are one atomic step.

\subsection{Replay}

Replay is addressed structurally in Section~\ref{sec:replay} and is
restated here as a security property: there is no sequence of proposer
re-builds, vertex re-issues, crash-before-commit, or network
re-delivery that credits a deposit twice or exports a withdraw twice,
because (i) the import consumes its UTXO exactly once, (ii) the venue
\texttt{seen:<txID>} returns the original outcome on re-relay, and (iii)
both facts are committed to the state root, so any validator that
replayed differently would fork the root and lose consensus. Replay
safety is thus not a best-effort property of the proxy; it is enforced
by the consensus check itself.

\subsection{Failed-export recoverability}
\label{sec:strand}

The only residual hazard is a \emph{stranded} custody transition: a
block whose proposer mutated the venue (and possibly began an export)
but which did not gain consensus, or whose export leg failed after the
venue debit. The governing requirement is that funds are always
\emph{unspent and recoverable, never burned}:

\begin{itemize}[leftmargin=2em,nosep]
\item \textbf{Failed export after debit.} If the venue debited
  \texttt{available} but the export leg failed, the realized amount MUST
  be returned to the account (full refund of the would-be-withdrawn
  amount), exactly as the order-submission failure path full-refunds
  escrow rather than stranding it. The withdraw then either does not
  occur (account made whole) or is re-attempted; it never half-occurs.
\item \textbf{Lost-poll block.} A block that loses the consensus poll
  does not commit its custody transitions to the canonical chain. The
  source UTXO of a deposit remains unconsumed on the canonical chain and
  is re-proposable. For a withdraw, the atomic discipline of the reorg
  requirement ensures the export UTXO is not produced on the canonical
  chain if the debit is not committed there --- so the funds remain on
  the venue, available, and recoverable.
\item \textbf{Never burned.} No failure path destroys units. A stranded
  transition leaves value either on the originating chain (deposit not
  finalized) or on the venue (withdraw not finalized), always spendable.
  This is the conservation invariant of Section~\ref{sec:conservation}
  holding even under failure: $I = A + L + F + E$ never loses a term to
  a black hole.
\end{itemize}

\subsection{Determinism as a security property}

Finally, the determinism that motivates the whole LP is itself a safety
property: because every validator applies byte-identical carried records
and reaches an identical state root, there is no state at which honest
validators disagree about how much value entered or left the system.
Disagreement about custody is, for a value-bearing venue, the most
dangerous failure mode; the carried-payload model removes it by
construction rather than detecting it after the fact.
